\subsection{Projektavimo apimtis ir tikslai}
\label{sec:projektavimo-apimtis}

Projektavimo tikslas buvo įgyvendinti tik tuos architektūros ir platformos elementus, kurie būtini darbo eksperimentams. Dėl šios priežasties emuliatorius nėra pilnas RISC-V suderinamumo modelis: jis nepalaiko kelių procesorių, pilnos įrenginių įvairovės, FPU, suspaustų instrukcijų plėtinio ar tikslaus visų mikroarchitektūrinių efektų. Vietoj to pasirinkta praktinė kryptis: pakankamai tiksliai emuliuoti RV64IMA, privilegijuotą vykdymą, Sv39 virtualią atmintį ir nedidelį įrenginių rinkinį, kad būtų galima paleisti realią programinės įrangos grandinę.

Darbe naudota konfigūracija atitinka šiuos pagrindinius projektavimo sprendimus:

\begin{itemize}
    \item \textbf{vienas aparatinis srautas} -- sistemoje emuliuojamas vienas hartas, todėl atsisakyta kelių procesorių sinchronizacijos ir sudėtingesnio pertraukimų paskirstymo;
    \item \textbf{RV64 kryptis} -- emuliatoriaus pagrindinė konfigūracija yra 64 bitų, nes Linux paleidimo ir vartotojo erdvės eksperimentai šiame darbe vykdomi RV64 aplinkoje;
    \item \textbf{minimalus Linux įrenginių rinkinys} -- emuliuojami tik tie įrenginiai, kurių reikia konsolės įvesčiai/išvesčiai, laikmačiui, išoriniams pertraukimams ir sistemos išjungimui;
    \item \textbf{eksperimentinis QPU plėtinys} -- specializuotos instrukcijos įtraukiamos į tą patį instrukcijų dekodavimo kelią kaip ir standartinės RISC-V instrukcijos, tačiau jų semantika deleguojama atskiram QPU būsenos valdymo komponentui;
    \item \textbf{atkuriama surinkimo aplinka} -- OpenSBI, Linux branduolys, DTB ir initramfs kuriami deklaratyviai, nes visi šie komponentai turi būti suderinti tarpusavyje.
\end{itemize}
